 \setlength{\baselineskip}{20pt}
\chapter{总结与展望}
\label{cha:chap5}



\section{研究总结}

随着智能驾驶技术的快速发展，自动泊车系统已成为提升驾驶便利性与安全性的核心功能之一。其中，停车位检测作为自动泊车的前提与基础，其性能直接影响整个系统的可靠性。然而，在复杂真实场景下，停车位检测面临光照变化、遮挡干扰、非标准车位等多重挑战，传统基于规则的方法已难以满足实际应用需求。本文围绕复杂场景下停车位检测这一关键技术问题，从数据构建、算法设计到系统优化展开了系统性研究，实现了停车位检测从基于规则的静态流程向结构感知且具备时序一致性的端到端检测方法的范式演进。具体工作内容如下：

（1）数据与学习范式创新

针对现有停车位检测研究中高质量标注数据匮乏、复杂场景覆盖不足的问题，本文构建了复杂真实场景停车位数据集 CRPS-D，涵盖地上停车场、地下停车库、路边停车区等多种场景，并覆盖白天、黑夜、晴天、阴天、雨天等多种天气与光照条件。该数据集规模超过现有公开数据集（如ps2.0 和 SNU），且包含更多斜向车位和复杂场景样本，为复杂场景下的停车位检测研究提供了统一的评估基准。
在此基础上，本文提出了面向停车位检测任务的半监督检测框架（SS-PSD），通过教师-学生学习范式有效利用未标注数据，缓解了标注成本高昂对模型训练与泛化能力的制约。该框架集成了置信度引导掩码（CGM）一致性约束和自适应特征扰动机制（Adaptive-VAT），教师模型通过学生模型参数的指数移动平均（EMA）生成，利用一致性正则化引导学生模型学习。为避免强制错误预测保持一致的问题，引入可训练的置信度图为不同区域分配权重，并对低置信度区域进行掩码处理，构建置信度引导的掩码一致性约束。同时，设计了Adaptive-VAT自适应特征扰动机制，根据教师与学生模型对扰动的鲁棒性，选择性生成更具挑战性的对抗噪声，促进学生模型学习更稳定的特征表示。

（2）空间结构感知建模

针对传统停车位检测方法中依赖启发式规则进行后处理匹配、结构一致性难以保证的问题，本文提出了面向结构感知的停车位关键点检测与匹配方法（SAKM-PSD）。该方法从建模角度将停车位检测问题重构为基于关键点的结构化预测问题，在端到端训练过程中显式引入停车位内部结构关系约束。
具体而言，该方法通过结构感知 CutMix 策略（SA-CutMix）增强背景多样性，在数据增强过程中保持停车位结构的完整性，避免传统 CutMix 可能破坏停车位几何关系的问题。通过注意力引导插值模块（AGI Module）提升特征对关键结构位置的敏感性，使模型能够更准确地捕捉停车位的边角点和边界线。通过稀配聚焦损失（SPF-Loss）缓解样本稀疏带来的训练不稳定问题，对不同难度的样本分配不同的权重，提高模型对稀疏关键点的学习能力。这些设计使结构一致性直接参与模型优化，从而避免了规则依赖并提升了复杂与密集场景下的检测稳定性与一致性。

（3）多帧时序特征融合

针对真实驾驶场景中停车位检测结果在连续帧条件下易受遮挡、视角变化和噪声影响而不稳定的问题，本文在单帧结构感知检测方法的基础上，进一步提出了多帧时序特征融合的停车位检测方法（MTF-PSD）。该方法通过对连续帧中的结构感知表示进行时序建模，将停车位检测从静态空间结构建模扩展至时空联合建模。具体实现上，该方法引入了跨时帧自注意力机制，用于在特征层面实现多帧时序信息的自适应融合。跨时帧自注意力模块以当前帧特征作为查询项，以前一帧特征作为键和值，通过注意力机制显式建模两帧之间的空间关联关系。同时设计了动态结构感知的时序融合模块（DSAF），通过联合建模全局通道重要性与局部结构显著性，自适应地调节当前帧特征与跨帧注意力特征之间的融合权重。动态融合模块从当前帧特征中提取全局通道权重，刻画不同语义通道在当前时刻的重要程度；同时预测结构感知掩码，突出停车位结构相关的空间区域。将通道级权重与空间结构掩码进行联合建模，得到最终的动态融合权重，使模型能够在空间和通道两个层面上自适应地平衡当前帧与历史帧信息。这些设计有效提升了检测结果在时间维度上的一致性与鲁棒性，使所提出的方法更加贴近实际应用需求。

综上所述，本文通过系统性的研究，逐步将停车位检测从基于规则的静态检测流程，发展为结构感知且具备时序一致性的端到端检测方法。实验结果表明，所提出的方法在多种复杂场景下均具备良好的检测精度与稳定性，为自动泊车系统的实际应用提供了有力支持。


\section{未来展望}

随着自动驾驶技术的快速发展，自动泊车系统作为其核心功能之一，对停车位检测的精度、鲁棒性和实时性提出了更高的要求。尽管本文在复杂场景下停车位检测技术研究方面取得了显著进展，通过构建大规模数据集、提出半监督学习框架、设计结构感知端到端检测方法以及多帧时序融合策略，有效提升了检测性能，但面对实际应用中的极端天气、复杂光照、非标准停车位等挑战，该领域仍有多个重要方向值得进一步探索与深化：

### \textbf{多模态融合与系统集成：}当前研究主要基于视觉传感器的单模态信息，未来可考虑融合激光雷达、毫米波雷达等多模态传感器数据，构建多源异构信息融合的停车位检测系统。通过互补不同传感器的优势，如激光雷达的高精度距离测量能力和毫米波雷达的全天候工作能力，有望进一步提升在极端天气、强遮挡等场景下的检测鲁棒性。同时，停车位检测仅是自动泊车系统的一个环节，未来可将检测算法与路径规划、运动控制等模块深度集成，构建端到端的自动泊车系统。通过联合优化感知、规划与控制模块，实现从环境感知到车辆控制的全流程优化，提升整体系统的性能与可靠性。例如，可将停车位检测结果直接作为路径规划的输入，减少模块间的信息损失，提高系统响应速度。


\textbf{模型轻量化与实时性优化：}随着自动驾驶技术的快速发展，车载计算资源的限制日益凸显。未来可针对停车位检测算法进行轻量化设计，通过模型压缩、知识蒸馏、量化等技术，在保证检测精度的同时显著提升推理速度，满足实际车载系统的实时性要求。例如，可采用结构化剪枝方法移除模型中冗余的神经元和连接，或通过知识蒸馏将大模型的知识迁移到小模型中，实现模型性能与效率的平衡。同时，可探索硬件加速方案，如利用GPU、FPGA等专用硬件实现停车位检测算法的并行计算，进一步提升处理速度。针对特定场景的快速检测需求，还可设计场景自适应的检测策略，在简单场景下使用轻量级模型，在复杂场景下切换到更强大的模型，实现计算资源的智能分配。


\textbf{学习范式与场景适应：}尽管本文提出了半监督学习框架，但标注成本仍然较高。未来可探索自监督与无监督学习方法，通过利用大量未标注数据自动学习停车位的结构特征，进一步降低对人工标注的依赖。例如，可设计基于对比学习的自监督方法，通过数据增强生成正样本对，让模型学习停车位的不变特征表示。
同时，当前模型在跨场景迁移时仍面临挑战，未来可引入域适应技术，使模型能够快速适应新的停车场景，减少对目标场景标注数据的依赖。例如，可采用对抗域适应方法，通过域鉴别器区分源域和目标域的特征分布，引导模型学习域不变的特征表示，提升模型在未知场景中的泛化能力。


\textbf{复杂场景与应用拓展：}未来可进一步拓展停车位检测的应用场景，如非标准停车位、临时停车位、立体停车场等复杂场景的检测与识别。针对这些场景的特点，需要设计专门的检测算法，例如针对立体停车场的多层停车位检测，需要考虑高度信息和视角变化的影响；针对临时停车位，需要处理不规则的边界和动态变化的场景。
此外，可将停车位检测技术与其他智能交通系统相结合，例如与智能 parking 引导系统集成，为驾驶员提供实时的停车位信息和导航服务；与自动代客泊车系统集成，实现从车辆 drop-off 到自动 parking 的全流程自动化。这些应用拓展将进一步推动停车位检测技术的实用化进程。

综上所述，停车位检测技术作为自动泊车系统的核心环节，其发展对于提升智能驾驶的安全性与便利性具有重要意义。未来的研究将继续围绕多模态融合、模型轻量化、学习范式创新与场景适应等方向展开，不断推动停车位检测技术向更智能、更可靠、更实用的方向发展，为智能交通系统的建设贡献力量。
